\IfFileExists{stacks-project.cls}{%
\documentclass{stacks-project}
}{%
\documentclass{amsart}
}

% For dealing with references we use the comment environment
\usepackage{verbatim}
\newenvironment{reference}{\comment}{\endcomment}
%\newenvironment{reference}{}{}
\newenvironment{slogan}{\comment}{\endcomment}
\newenvironment{history}{\comment}{\endcomment}

% For commutative diagrams we use Xy-pic
\usepackage[all]{xy}

% We use 2cell for 2-commutative diagrams.
\xyoption{2cell}
\UseAllTwocells

% We use multicol for the list of chapters between chapters
\usepackage{multicol}

% This is generally recommended for better output
\usepackage{lmodern}
\usepackage[T1]{fontenc}

% For cross-file-references
\usepackage{xr-hyper}

% Package for hypertext links:
\usepackage{hyperref}

% For any local file, say "hello.tex" you want to link to please
% use \externaldocument[hello-]{hello}
\externaldocument[introduction-]{introduction}
\externaldocument[conventions-]{conventions}
\externaldocument[sets-]{sets}
\externaldocument[categories-]{categories}
\externaldocument[topology-]{topology}
\externaldocument[sheaves-]{sheaves}
\externaldocument[sites-]{sites}
\externaldocument[stacks-]{stacks}
\externaldocument[fields-]{fields}
\externaldocument[algebra-]{algebra}
\externaldocument[brauer-]{brauer}
\externaldocument[homology-]{homology}
\externaldocument[derived-]{derived}
\externaldocument[simplicial-]{simplicial}
\externaldocument[more-algebra-]{more-algebra}
\externaldocument[smoothing-]{smoothing}
\externaldocument[modules-]{modules}
\externaldocument[sites-modules-]{sites-modules}
\externaldocument[injectives-]{injectives}
\externaldocument[cohomology-]{cohomology}
\externaldocument[sites-cohomology-]{sites-cohomology}
\externaldocument[dga-]{dga}
\externaldocument[dpa-]{dpa}
\externaldocument[sdga-]{sdga}
\externaldocument[hypercovering-]{hypercovering}
\externaldocument[schemes-]{schemes}
\externaldocument[constructions-]{constructions}
\externaldocument[properties-]{properties}
\externaldocument[morphisms-]{morphisms}
\externaldocument[coherent-]{coherent}
\externaldocument[divisors-]{divisors}
\externaldocument[limits-]{limits}
\externaldocument[varieties-]{varieties}
\externaldocument[topologies-]{topologies}
\externaldocument[descent-]{descent}
\externaldocument[perfect-]{perfect}
\externaldocument[more-morphisms-]{more-morphisms}
\externaldocument[flat-]{flat}
\externaldocument[groupoids-]{groupoids}
\externaldocument[more-groupoids-]{more-groupoids}
\externaldocument[etale-]{etale}
\externaldocument[chow-]{chow}
\externaldocument[intersection-]{intersection}
\externaldocument[pic-]{pic}
\externaldocument[weil-]{weil}
\externaldocument[adequate-]{adequate}
\externaldocument[dualizing-]{dualizing}
\externaldocument[duality-]{duality}
\externaldocument[discriminant-]{discriminant}
\externaldocument[derham-]{derham}
\externaldocument[local-cohomology-]{local-cohomology}
\externaldocument[algebraization-]{algebraization}
\externaldocument[curves-]{curves}
\externaldocument[resolve-]{resolve}
\externaldocument[models-]{models}
\externaldocument[functors-]{functors}
\externaldocument[equiv-]{equiv}
\externaldocument[pione-]{pione}
\externaldocument[etale-cohomology-]{etale-cohomology}
\externaldocument[proetale-]{proetale}
\externaldocument[relative-cycles-]{relative-cycles}
\externaldocument[more-etale-]{more-etale}
\externaldocument[trace-]{trace}
\externaldocument[crystalline-]{crystalline}
\externaldocument[spaces-]{spaces}
\externaldocument[spaces-properties-]{spaces-properties}
\externaldocument[spaces-morphisms-]{spaces-morphisms}
\externaldocument[decent-spaces-]{decent-spaces}
\externaldocument[spaces-cohomology-]{spaces-cohomology}
\externaldocument[spaces-limits-]{spaces-limits}
\externaldocument[spaces-divisors-]{spaces-divisors}
\externaldocument[spaces-over-fields-]{spaces-over-fields}
\externaldocument[spaces-topologies-]{spaces-topologies}
\externaldocument[spaces-descent-]{spaces-descent}
\externaldocument[spaces-perfect-]{spaces-perfect}
\externaldocument[spaces-more-morphisms-]{spaces-more-morphisms}
\externaldocument[spaces-flat-]{spaces-flat}
\externaldocument[spaces-groupoids-]{spaces-groupoids}
\externaldocument[spaces-more-groupoids-]{spaces-more-groupoids}
\externaldocument[bootstrap-]{bootstrap}
\externaldocument[spaces-pushouts-]{spaces-pushouts}
\externaldocument[spaces-chow-]{spaces-chow}
\externaldocument[groupoids-quotients-]{groupoids-quotients}
\externaldocument[spaces-more-cohomology-]{spaces-more-cohomology}
\externaldocument[spaces-simplicial-]{spaces-simplicial}
\externaldocument[spaces-duality-]{spaces-duality}
\externaldocument[formal-spaces-]{formal-spaces}
\externaldocument[restricted-]{restricted}
\externaldocument[spaces-resolve-]{spaces-resolve}
\externaldocument[formal-defos-]{formal-defos}
\externaldocument[defos-]{defos}
\externaldocument[cotangent-]{cotangent}
\externaldocument[examples-defos-]{examples-defos}
\externaldocument[algebraic-]{algebraic}
\externaldocument[examples-stacks-]{examples-stacks}
\externaldocument[stacks-sheaves-]{stacks-sheaves}
\externaldocument[criteria-]{criteria}
\externaldocument[artin-]{artin}
\externaldocument[quot-]{quot}
\externaldocument[stacks-properties-]{stacks-properties}
\externaldocument[stacks-morphisms-]{stacks-morphisms}
\externaldocument[stacks-limits-]{stacks-limits}
\externaldocument[stacks-cohomology-]{stacks-cohomology}
\externaldocument[stacks-perfect-]{stacks-perfect}
\externaldocument[stacks-introduction-]{stacks-introduction}
\externaldocument[stacks-more-morphisms-]{stacks-more-morphisms}
\externaldocument[stacks-geometry-]{stacks-geometry}
\externaldocument[moduli-]{moduli}
\externaldocument[moduli-curves-]{moduli-curves}
\externaldocument[examples-]{examples}
\externaldocument[exercises-]{exercises}
\externaldocument[guide-]{guide}
\externaldocument[desirables-]{desirables}
\externaldocument[coding-]{coding}
\externaldocument[obsolete-]{obsolete}
\externaldocument[fdl-]{fdl}
\externaldocument[index-]{index}

% Theorem environments.
%
\theoremstyle{plain}
\newtheorem{theorem}[subsection]{Theorem}
\newtheorem{proposition}[subsection]{Proposition}
\newtheorem{lemma}[subsection]{Lemma}

\theoremstyle{definition}
\newtheorem{definition}[subsection]{Definition}
\newtheorem{example}[subsection]{Example}
\newtheorem{exercise}[subsection]{Exercise}
\newtheorem{situation}[subsection]{Situation}

\theoremstyle{remark}
\newtheorem{remark}[subsection]{Remark}
\newtheorem{remarks}[subsection]{Remarks}

\numberwithin{equation}{subsection}

% Macros
%
\def\lim{\mathop{\mathrm{lim}}\nolimits}
\def\colim{\mathop{\mathrm{colim}}\nolimits}
\def\Spec{\mathop{\mathrm{Spec}}}
\def\Hom{\mathop{\mathrm{Hom}}\nolimits}
\def\Ext{\mathop{\mathrm{Ext}}\nolimits}
\def\SheafHom{\mathop{\mathcal{H}\!\mathit{om}}\nolimits}
\def\SheafExt{\mathop{\mathcal{E}\!\mathit{xt}}\nolimits}
\def\Sch{\mathit{Sch}}
\def\Mor{\mathop{\mathrm{Mor}}\nolimits}
\def\Ob{\mathop{\mathrm{Ob}}\nolimits}
\def\Sh{\mathop{\mathit{Sh}}\nolimits}
\def\NL{\mathop{N\!L}\nolimits}
\def\CH{\mathop{\mathrm{CH}}\nolimits}
\def\proetale{{pro\text{-}\acute{e}tale}}
\def\etale{{\acute{e}tale}}
\def\QCoh{\mathit{QCoh}}
\def\Ker{\mathop{\mathrm{Ker}}}
\def\Im{\mathop{\mathrm{Im}}}
\def\Coker{\mathop{\mathrm{Coker}}}
\def\Coim{\mathop{\mathrm{Coim}}}

% Boxtimes
%
\DeclareMathSymbol{\boxtimes}{\mathbin}{AMSa}{"02}

%
% Macros for moduli stacks/spaces
%
\def\QCohstack{\mathcal{QC}\!\mathit{oh}}
\def\Cohstack{\mathcal{C}\!\mathit{oh}}
\def\Spacesstack{\mathcal{S}\!\mathit{paces}}
\def\Quotfunctor{\mathrm{Quot}}
\def\Hilbfunctor{\mathrm{Hilb}}
\def\Curvesstack{\mathcal{C}\!\mathit{urves}}
\def\Polarizedstack{\mathcal{P}\!\mathit{olarized}}
\def\Complexesstack{\mathcal{C}\!\mathit{omplexes}}
% \Pic is the operator that assigns to X its picard group, usage \Pic(X)
% \Picardstack_{X/B} denotes the Picard stack of X over B
% \Picardfunctor_{X/B} denotes the Picard functor of X over B
\def\Pic{\mathop{\mathrm{Pic}}\nolimits}
\def\Picardstack{\mathcal{P}\!\mathit{ic}}
\def\Picardfunctor{\mathrm{Pic}}
\def\Deformationcategory{\mathcal{D}\!\mathit{ef}}
\usepackage{cleveref}
\crefname{originallemma}{original lemma}{original lemmas}
\Crefname{originallemma}{Original lemma}{Original lemmas}
\makeatletter
\newlabel{algebra-lemma-directed-colimit-exact}{{00DB}{source}{}{sourcepointer-1}{}}
\newlabel{algebra-lemma-directed-colimit-exact@cref}{{[originallemma][0][]00DB}{[1][0][]source}{}{}{}}
\makeatother
\crefname{originallemma}{original lemma}{original lemmas}
\Crefname{originallemma}{Original lemma}{Original lemmas}
\makeatletter
\newlabel{derived-lemma-find-existence-computes}{{06XN}{source}{}{sourcepointer-2}{}}
\newlabel{derived-lemma-find-existence-computes@cref}{{[originallemma][0][]06XN}{[1][0][]source}{}{}{}}
\makeatother
\crefname{originallemma}{original lemma}{original lemmas}
\Crefname{originallemma}{Original lemma}{Original lemmas}
\makeatletter
\newlabel{homology-lemma-good-resolution-gives-qis}{{09IZ}{source}{}{sourcepointer-3}{}}
\newlabel{homology-lemma-good-resolution-gives-qis@cref}{{[originallemma][0][]09IZ}{[1][0][]source}{}{}{}}
\makeatother
\crefname{originallemma}{original lemma}{original lemmas}
\Crefname{originallemma}{Original lemma}{Original lemmas}
\makeatletter
\newlabel{derived-lemma-representable-homological}{{0149}{source}{}{sourcepointer-4}{}}
\newlabel{derived-lemma-representable-homological@cref}{{[originallemma][0][]0149}{[1][0][]source}{}{}{}}
\makeatother
\crefname{originallemma}{original lemma}{original lemmas}
\Crefname{originallemma}{Original lemma}{Original lemmas}
\makeatletter
\newlabel{homology-lemma-good-right-resolution-gives-qis}{{09J0}{source}{}{sourcepointer-5}{}}
\newlabel{homology-lemma-good-right-resolution-gives-qis@cref}{{[originallemma][0][]09J0}{[1][0][]source}{}{}{}}
\makeatother
\crefname{originallemma}{original lemma}{original lemmas}
\Crefname{originallemma}{Original lemma}{Original lemmas}
\makeatletter
\newlabel{derived-lemma-easier-axiom-four}{{014C}{source}{}{sourcepointer-6}{}}
\newlabel{derived-lemma-easier-axiom-four@cref}{{[originallemma][0][]014C}{[1][0][]source}{}{}{}}
\makeatother
\crefname{originallemma}{original lemma}{original lemmas}
\Crefname{originallemma}{Original lemma}{Original lemmas}
\makeatletter
\newlabel{derived-lemma-make-injective}{{013N}{source}{}{sourcepointer-7}{}}
\newlabel{derived-lemma-make-injective@cref}{{[originallemma][0][]013N}{[1][0][]source}{}{}{}}
\makeatother
\title{The composition of two derived tensor functors is a derived tensor functor}
\author{The Stacks Project Authors}
\date{stacks-project:dga.tex:local-snapshot-20261012; GNU FDL 1.2 or later}
\begin{document}
\maketitle
\noindent\textit{Editorial note.} The counted statement is the author's own lemma environment labelled lemma-compose-tensor-functors-general-algebra (source lines 5958-5975); the proof environment begins on the second line after it and is its own proof, closing at source line 6064.\par\medskip
\noindent\textit{Editorial note.} The statement and the proof invoke six author results of the same chapter by label; all six are retained as context spans and the generator expands the context transitively (budget 800 lines for this card), while reference targets in other chapters that carry canonical project tags are auto-declared as non-load-bearing source pointers.\par\medskip
\noindent\textit{Editorial note.} The retained preamble is the chapter's own minimal preamble (sources/stacks-dga/preamble.tex, taken with its real line count), a verbatim excerpt of the native Xy-pic setup, the native theorem declarations with the mathematical macros, and the project's own no-op reference environments; no macro was invented or redefined.\par\medskip
\noindent\textit{Editorial scope.} One original result of the Stacks Project's Differential Graded Algebra chapter, counted once: the composition of two derived tensor functors is isomorphic to a derived tensor functor given by the tensor product of the two bimodules.\par\medskip
\noindent\textit{Allowed standard prerequisites (editorial).}\begin{itemize}
\item Differential graded algebras and bimodules. Dg R-algebras (A,d), (B,d), (C,d), and dg (A,B)- and (B,C)-bimodules, with their homotopy and derived categories D(-,d).. It is the setting of the statement (source lines 5960-5964).
\item Derived tensor product of bimodules. The left derived functor - tensor\textasciicircum{}L\_A N on D(A,d) defined by the construction of the chapter, together with its functoriality in the bimodule up to quasi-isomorphism.. The two functors being composed and the announced N'' are derived tensor functors (source lines 5965-5974, 5978-5982).
\item Property (P) for bimodules and its filtration sequences. Every dg bimodule has a quasi-isomorphic representative with property (P), carrying an admissible filtration whose graded pieces are direct sums of shifts of the free bimodule, and a split admissible short exact sequence displaying the filtration (the author's lemmas lemma-bimodule-resolve and lemma-bimodule-property-P-sequence, retained as context).. Both reduction chains of the proof run along these filtrations (source lines 5980-5982, 5992-6024, 6035-6051).
\item Exactness of the derived tensor functor. The derived tensor functor is exact as a functor of triangulated categories, commutes with direct sums, and transforms admissible exact sequences of complexes into admissible exact sequences (the author's lemma-compose-tensor-functors-general, retained as context).. It is what licenses the two reductions (source lines 5988-6018).
\end{itemize}\medskip
\ifdefined\hypertarget\hypertarget{sourcepointer-1}{}\fi
\noindent\textit{Source reference.} Original lemma 00DB: sources/stacks-chow/tags.txt, source lines 491--491. Cross-chapter pointer whose target carries the canonical tag 00DB in the project tag list; declared so the reference resolves inside the assembled package, while every step the counted argument uses remains inside the retained spans.\par\medskip
\ifdefined\hypertarget\hypertarget{sourcepointer-2}{}\fi
\noindent\textit{Source reference.} Original lemma 06XN: sources/stacks-chow/tags.txt, source lines 8857--8857. Cross-chapter pointer whose target carries the canonical tag 06XN in the project tag list; declared so the reference resolves inside the assembled package, while every step the counted argument uses remains inside the retained spans.\par\medskip
\ifdefined\hypertarget\hypertarget{sourcepointer-3}{}\fi
\noindent\textit{Source reference.} Original lemma 09IZ: sources/stacks-chow/tags.txt, source lines 12057--12057. Cross-chapter pointer whose target carries the canonical tag 09IZ in the project tag list; declared so the reference resolves inside the assembled package, while every step the counted argument uses remains inside the retained spans.\par\medskip
\ifdefined\hypertarget\hypertarget{sourcepointer-4}{}\fi
\noindent\textit{Source reference.} Original lemma 0149: sources/stacks-chow/tags.txt, source lines 1464--1464. Cross-chapter pointer whose target carries the canonical tag 0149 in the project tag list; declared so the reference resolves inside the assembled package, while every step the counted argument uses remains inside the retained spans.\par\medskip
\ifdefined\hypertarget\hypertarget{sourcepointer-5}{}\fi
\noindent\textit{Source reference.} Original lemma 09J0: sources/stacks-chow/tags.txt, source lines 12058--12058. Cross-chapter pointer whose target carries the canonical tag 09J0 in the project tag list; declared so the reference resolves inside the assembled package, while every step the counted argument uses remains inside the retained spans.\par\medskip
\ifdefined\hypertarget\hypertarget{sourcepointer-6}{}\fi
\noindent\textit{Source reference.} Original lemma 014C: sources/stacks-chow/tags.txt, source lines 1467--1467. Cross-chapter pointer whose target carries the canonical tag 014C in the project tag list; declared so the reference resolves inside the assembled package, while every step the counted argument uses remains inside the retained spans.\par\medskip
\ifdefined\hypertarget\hypertarget{sourcepointer-7}{}\fi
\noindent\textit{Source reference.} Original lemma 013N: sources/stacks-chow/tags.txt, source lines 1442--1442. Cross-chapter pointer whose target carries the canonical tag 013N in the project tag list; declared so the reference resolves inside the assembled package, while every step the counted argument uses remains inside the retained spans.\par\medskip
\begin{equation}
\label{equation-plain-versus-derived-algebras}
(A \otimes_R B)_B \otimes_B^\mathbf{L} {}_B(B \otimes_R C)_C
\longrightarrow
(A \otimes_R B \otimes_R C)_C
\end{equation}

\begin{lemma}
\label{lemma-functoriality-bc}
Let $R$ be a ring. Let $(A, \text{d})$ and $(B, \text{d})$ be
differential graded $R$-algebras. Let $f : N \to N'$ be a
homomorphism of differential graded $(A, B)$-bimodules.
Then $f$ induces a morphism of functors
$$
1\otimes f :
- \otimes_A^\mathbf{L} N
\longrightarrow
- \otimes_A^\mathbf{L} N'
$$
If $f$ is a quasi-isomorphism, then $1 \otimes f$ is an isomorphism of
functors.
\end{lemma}

\begin{proof}
Let $M$ be a differential graded $A$-module with property (P).
Then $1 \otimes f : M \otimes_A N \to M \otimes_A N'$ is a
map of differential graded $B$-modules. Moreover, this is functorial
with respect to $M$. Since the functors
$- \otimes_A^\mathbf{L} N$ and $- \otimes_A^\mathbf{L} N'$ are
computed by tensoring on objects with property (P)
(Lemma \ref{lemma-derived-bc}) we obtain a transformation of functors
as indicated.

\medskip\noindent
Assume that $f$ is a quasi-isomorphism. Let $F_\bullet$ be the
given filtration on $M$. Observe that
$M \otimes_A N = \colim F_i(M) \otimes_A N$ and
$M \otimes_A N' = \colim F_i(M) \otimes_A N'$.
Hence it suffices to show that
$F_n(M) \otimes_A N \to F_n(M) \otimes_A N'$
is a quasi-isomorphism (filtered colimits are exact, see
Algebra, Lemma \ref{algebra-lemma-directed-colimit-exact}).
Since the inclusions $F_n(M) \to F_{n + 1}(M)$
are split as maps of graded $A$-modules we see that
$$
0 \to F_n(M) \otimes_A N \to F_{n + 1}(M) \otimes_A N \to
F_{n + 1}(M)/F_n(M) \otimes_A N \to 0
$$
is a short exact sequence of differential graded $B$-modules.
There is a similar sequence for $N'$ and $f$ induces a map
of short exact sequences. Hence by induction on $n$ (starting with $n = -1$
when $F_{-1}(M) = 0$) we conclude that it suffices to show that
the map $F_{n + 1}(M)/F_n(M) \otimes_A N \to F_{n + 1}(M)/F_n(M) \otimes_A N'$
is a quasi-isomorphism. This is true because $F_{n + 1}(M)/F_n(M)$
is a direct sum of shifts of $A$ and the result is true for $A[k]$
as $f : N \to N'$ is a quasi-isomorphism.
\end{proof}

\begin{lemma}
\label{lemma-bimodule-resolve}
Let $R$ be a ring. Let $(A, \text{d})$ and $(B, \text{d})$ be
differential graded $R$-algebras. Let $M$ be a differential graded
$(A, B)$-bimodule. There exists a homomorphism $P \to M$
of differential graded $(A, B)$-bimodules which is a quasi-isomorphism
such that $P$ has property (P) as defined above.
\end{lemma}

\begin{proof}
Immediate from Lemmas \ref{lemma-bimodule-over-tensor} and
\ref{lemma-resolve}.
\end{proof}

\begin{lemma}
\label{lemma-compose-tensor-functors-general}
Let $R$ be a ring. Let $(A, \text{d})$, $(B, \text{d})$, and
$(C, \text{d})$ be differential graded $R$-algebras.
Let $N$ be a differential graded $(A, B)$-bimodule.
Let $N'$ be a differential graded $(B, C)$-module.
Assume (\ref{equation-plain-versus-derived}) is an isomorphism.
Then the composition
$$
\xymatrix{
D(A, \text{d}) \ar[rr]^{- \otimes_A^\mathbf{L} N} & &
D(B, \text{d}) \ar[rr]^{- \otimes_B^\mathbf{L} N'} & &
D(C, \text{d})
}
$$
is isomorphic to $- \otimes_A^\mathbf{L} N''$ with
$N'' = N \otimes_B N'$ viewed as $(A, C)$-bimodule.
\end{lemma}

\begin{proof}
Let us define a transformation of functors
$$
(- \otimes_A^\mathbf{L} N) \otimes_B^\mathbf{L} N'
\longrightarrow
- \otimes_A^\mathbf{L} N''
$$
To do this, let
$M$ be a differential graded $A$-module with property (P).
According to the construction of the functor $- \otimes_A^\mathbf{L} N''$
of the proof of Lemma \ref{lemma-derived-bc} the plain tensor
product $M \otimes_A N''$ represents $M \otimes_A^\mathbf{L} N''$
in $D(C, \text{d})$. Then we write
$$
M \otimes_A N'' =
M \otimes_A (N \otimes_B N') =
(M \otimes_A N) \otimes_B N'
$$
The module $M \otimes_A N$ represents $M \otimes_A^\mathbf{L} N$
in $D(B, \text{d})$. Choose a quasi-isomorphism $Q \to M \otimes_A N$
where $Q$ is a differential graded $B$-module with property (P). Then
$Q \otimes_B N'$ represents
$(M \otimes_A^\mathbf{L} N) \otimes_B^\mathbf{L} N'$ in
$D(C, \text{d})$.
Thus we can define our map via
$$
(M \otimes_A^\mathbf{L} N) \otimes_B^\mathbf{L} N' =
Q \otimes_B N' \to
M \otimes_A N \otimes_B N' =
M \otimes_A^\mathbf{L} N''
$$
The construction of this map is functorial in $M$ and compatible
with distinguished triangles and direct sums; we omit the details.
Consider the property $T$ of objects $M$ of $D(A, \text{d})$ 
expressing that this map is an isomorphism. Then
\begin{enumerate}
\item if $T$ holds for $M_i$ then $T$ holds for $\bigoplus M_i$,
\item if $T$ holds for $2$-out-of-$3$ in a distinguished
triangle, then it holds for the third, and
\item $T$ holds for $A[k]$ because here we obtain a shift
of the map (\ref{equation-plain-versus-derived}) which we
have assumed is an isomorphism.
\end{enumerate}
Thus by Remark \ref{remark-P-resolution} property $T$
always holds and the proof is complete.
\end{proof}

\begin{lemma}
\label{lemma-bimodule-property-P-sequence}
Let $R$ be a ring. Let $(A, \text{d})$ and $(B, \text{d})$ be
differential graded $R$-algebras. Let $P$ be a
differential graded $(A, B)$-bimodule having property (P)
with corresponding filtration $F_\bullet$, then we obtain a
short exact sequence
$$
0 \to
\bigoplus\nolimits F_iP \to
\bigoplus\nolimits F_iP \to P \to 0
$$
of differential graded $(A, B)$-bimodules which is split as a sequence
of graded $(A, B)$-bimodules.
\end{lemma}

\begin{proof}
Immediate from Lemmas \ref{lemma-bimodule-over-tensor} and
\ref{lemma-property-P-sequence}.
\end{proof}

\begin{lemma}
\label{lemma-derived-bc}
In the situation above, the left derived functor of $F$ exists.
We denote it
$- \otimes_A^\mathbf{L} N : D(A, \text{d}) \to D(B, \text{d})$.
\end{lemma}

\begin{proof}
We will use
Derived Categories, Lemma \ref{derived-lemma-find-existence-computes}
to prove this. As our collection $\mathcal{P}$
of objects we will use the objects with property (P).
Property (1) was shown in Lemma \ref{lemma-resolve}.
Property (2) holds because if $s : P \to P'$ is a quasi-isomorphism
of modules with property (P), then $s$ is a homotopy equivalence
by Lemma \ref{lemma-hom-derived}.
\end{proof}

\begin{lemma}
\label{lemma-bimodule-over-tensor}
Let $R$ be a ring. Let $(A, \text{d})$ and $(B, \text{d})$
be differential graded algebras over $R$. The construction above
defines an equivalence of categories
$$
\begin{matrix}
\text{differential graded}\\
(A, B)\text{-bimodules}
\end{matrix}
\longleftrightarrow
\begin{matrix}
\text{right differential graded }\\
A^{opp} \otimes_R B\text{-modules}
\end{matrix}
$$
\end{lemma}

\begin{proof}
Immediate from discussion the above.
\end{proof}

\begin{lemma}
\label{lemma-resolve}
Let $(A, \text{d})$ be a differential graded algebra.
Let $M$ be a differential graded $A$-module. There exists a homomorphism
$P \to M$ of differential graded $A$-modules such that
\begin{enumerate}
\item $P \to M$ is a quasi-isomorphism, and
\item $P$ has property (P).
\end{enumerate}
\end{lemma}

\begin{proof}
Set $M = M_0$. We inductively choose short exact sequences
$$
0 \to M_{i + 1} \to P_i \to M_i \to 0
$$
where the maps $P_i \to M_i$ are chosen as in Lemma \ref{lemma-good-quotient}.
This gives a ``resolution''
$$
\ldots \to P_2 \xrightarrow{f_2} P_1 \xrightarrow{f_1} P_0 \to M \to 0
$$
Then we set
$$
P = \bigoplus\nolimits_{i \geq 0} P_i
$$
as an $A$-module with grading given by
$P^n = \bigoplus_{a + b = n} P_{-a}^b$ and
differential (as in the construction of the total complex associated
to a double complex) by
$$
\text{d}_P(x) = f_{-a}(x) + (-1)^a \text{d}_{P_{-a}}(x)
$$
for $x \in P_{-a}^b$. With these conventions $P$ is indeed a differential
graded $A$-module. Recalling that each $P_i$ has a two step filtration
$0 \to P_i' \to P_i \to P_i'' \to 0$ we set
$$
F_{2i}P = \bigoplus\nolimits_{i \geq j \geq 0} P_j
\subset
\bigoplus\nolimits_{i \geq 0} P_i = P
$$
and we add $P'_{i + 1}$ to $F_{2i}P$ to get $F_{2i + 1}$.
These are differential graded submodules and the successive quotients
are direct sums of shifts of $A$. By
Lemma \ref{lemma-target-graded-projective} we see that
the inclusions $F_iP \to F_{i + 1}P$ are admissible monomorphisms.
Finally, we have to show that the map $P \to M$ (given by the
augmentation $P_0 \to M$) is a quasi-isomorphism. This follows from
Homology, Lemma \ref{homology-lemma-good-resolution-gives-qis}.
\end{proof}

\begin{equation}
\label{equation-plain-versus-derived}
N_B \otimes_B^\mathbf{L} N'
\longrightarrow
(N \otimes_B N')_C
\end{equation}

\begin{remark}
\label{remark-P-resolution}
Let $R$ be a ring. Let $(A, \text{d})$ be a differential graded $R$-algebra.
Using P-resolutions we can sometimes reduce statements about general
objects of $D(A, \text{d})$ to statements about $A[k]$. Namely, let
$T$ be a property of objects of $D(A, \text{d})$ and assume that
\begin{enumerate}
\item if $K_i$, $i \in I$ is a family of objects of $D(A, \text{d})$
and $T(K_i)$ holds for all $i \in I$, then $T(\bigoplus K_i)$,
\item if $K \to L \to M \to K[1]$ is a distinguished triangle of
$D(A, \text{d})$ and $T$ holds for two, then $T$
holds for the third object, and
\item $T(A[k])$ holds for all $k \in \mathbf{Z}$.
\end{enumerate}
Then $T$ holds for all objects of $D(A, \text{d})$. This is clear from
Lemmas \ref{lemma-property-P-sequence} and \ref{lemma-resolve}.
\end{remark}

\begin{lemma}
\label{lemma-property-P-sequence}
Let $(A, \text{d})$ be a differential graded algebra.
Let $P$ be a differential graded $A$-module. If $F_\bullet$
is a filtration as in property (P), then we obtain an
admissible short exact sequence
$$
0 \to
\bigoplus\nolimits F_iP \to
\bigoplus\nolimits F_iP \to P \to 0
$$
of differential graded $A$-modules.
\end{lemma}

\begin{proof}
The second map is the direct sum of the inclusion maps.
The first map on the summand $F_iP$ of the source is the sum
of the identity $F_iP \to F_iP$ and the negative of the inclusion
map $F_iP \to F_{i + 1}P$. Choose homomorphisms $s_i : F_{i + 1}P \to F_iP$
of graded $A$-modules which are left inverse to the inclusion maps.
Composing gives maps $s_{j, i} : F_jP \to F_iP$ for all $j > i$.
Then a left inverse of the first arrow maps $x \in F_jP$ to
$(s_{j, 0}(x), s_{j, 1}(x), \ldots, s_{j, j - 1}(x), 0, \ldots)$
in $\bigoplus F_iP$.
\end{proof}

\begin{lemma}
\label{lemma-hom-derived}
Let $(A, \text{d})$ be a differential graded algebra.
Let $M$ and $N$ be differential graded $A$-modules.
\begin{enumerate}
\item Let $P \to M$ be a P-resolution as in
Lemma \ref{lemma-resolve}. Then
$$
\Hom_{D(A, \text{d})}(M, N) =
\Hom_{K(\text{Mod}_{(A, \text{d})})}(P, N)
$$
\item Let $N \to I$ be an I-resolution as in
Lemma \ref{lemma-right-resolution}. Then
$$
\Hom_{D(A, \text{d})}(M, N) =
\Hom_{K(\text{Mod}_{(A, \text{d})})}(M, I)
$$
\end{enumerate}
\end{lemma}

\begin{proof}
Let $P \to M$ be as in (1). Since $P \to M$ is a quasi-isomorphism we see that
$$
\Hom_{D(A, \text{d})}(P, N) = \Hom_{D(A, \text{d})}(M, N)
$$
by definition of the derived category. A morphism
$f : P \to N$ in $D(A, \text{d})$ is equal to
$s^{-1}f'$ where $f' : P \to N'$ is a morphism and
$s : N \to N'$ is a quasi-isomorphism. Choose a distinguished triangle
$$
N \to N' \to Q \to N[1]
$$
As $s$ is a quasi-isomorphism, we see that $Q$ is acyclic. Thus
$\Hom_{K(\text{Mod}_{(A, \text{d})})}(P, Q[k]) = 0$ for all $k$ by
Lemma \ref{lemma-property-P-K-projective}. Since
$\Hom_{K(\text{Mod}_{(A, \text{d})})}(P, -)$
is cohomological, we conclude that we can lift $f' : P \to N'$
uniquely to a morphism $f : P \to N$. This finishes the proof.

\medskip\noindent
The proof of (2) is dual to that of (1) using
Lemma \ref{lemma-property-I-K-injective} in stead of
Lemma \ref{lemma-property-P-K-projective}.
\end{proof}

\begin{lemma}
\label{lemma-good-quotient}
Let $(A, \text{d})$ be a differential graded algebra.
Let $M$ be a differential graded $A$-module. There exists a homomorphism
$P \to M$ of differential graded $A$-modules with the following
properties
\begin{enumerate}
\item $P \to M$ is surjective,
\item $\Ker(\text{d}_P) \to \Ker(\text{d}_M)$ is surjective, and
\item $P$ sits in an admissible short exact sequence
$0 \to P' \to P \to P'' \to 0$ where $P'$, $P''$ are direct sums
of shifts of $A$.
\end{enumerate}
\end{lemma}

\begin{proof}
Let $P_k$ be the free $A$-module with generators $x, y$ in degrees
$k$ and $k + 1$. Define the structure of a differential graded
$A$-module on $P_k$ by setting $\text{d}(x) = y$ and $\text{d}(y) = 0$.
For every element $m \in M^k$ there is a homomorphism
$P_k \to M$ sending $x$ to $m$ and $y$ to $\text{d}(m)$.
Thus we see that there is a surjection from a direct sum
of copies of $P_k$ to $M$. This clearly produces $P \to M$
having properties (1) and (3). To obtain property (2) note
that if $m \in \Ker(\text{d}_M)$ has degree $k$, then there is a map
$A[k] \to M$ mapping $1$ to $m$. Hence we can achieve (2) by adding
a direct sum of copies of shifts of $A$.
\end{proof}

\begin{lemma}
\label{lemma-target-graded-projective}
Let $(A, \text{d})$ be a differential graded algebra.
Let $M \to P$ be a surjective homomorphism of differential graded
$A$-modules. If $P$ is projective as a graded $A$-module, then
$M \to P$ is an admissible epimorphism.
\end{lemma}

\begin{proof}
This is immediate from the definitions.
\end{proof}

\begin{lemma}
\label{lemma-property-I-K-injective}
Let $(A, \text{d})$ be a differential graded algebra.
Let $I$ be a differential graded $A$-module with property (I).
Then
$$
\Hom_{K(\text{Mod}_{(A, \text{d})})}(N, I) = 0
$$
for all acyclic differential graded $A$-modules $N$.
\end{lemma}

\begin{proof}
We will use that $K(\text{Mod}_{(A, \text{d})})$ is a triangulated
category (Proposition \ref{proposition-homotopy-category-triangulated}).
Let $F_\bullet$ be a filtration on $I$ as in property (I).
The short exact sequence of Lemma \ref{lemma-property-I-sequence}
produces a distinguished triangle. Hence by
Derived Categories, Lemma \ref{derived-lemma-representable-homological}
it suffices to show that
$$
\Hom_{K(\text{Mod}_{(A, \text{d})})}(N, I/F_iI) = 0
$$
for all acyclic differential graded $A$-modules $N$ and all $i$.
Each of the differential graded modules $I/F_iI$ has a finite filtration
by admissible monomorphisms, whose graded pieces are
products of $A^\vee[k]$. Thus it suffices to prove that
$$
\Hom_{K(\text{Mod}_{(A, \text{d})})}(N, A^\vee[k]) = 0
$$
for all acyclic differential graded $A$-modules $N$ and all $k$.
This follows from Lemma \ref{lemma-hom-into-shift-dual-free}
and the fact that $(-)^\vee$ is an exact functor.
\end{proof}

\begin{lemma}
\label{lemma-property-P-K-projective}
Let $(A, \text{d})$ be a differential graded algebra.
Let $P$ be a differential graded $A$-module with property (P).
Then
$$
\Hom_{K(\text{Mod}_{(A, \text{d})})}(P, N) = 0
$$
for all acyclic differential graded $A$-modules $N$.
\end{lemma}

\begin{proof}
We will use that $K(\text{Mod}_{(A, \text{d})})$ is a triangulated
category (Proposition \ref{proposition-homotopy-category-triangulated}).
Let $F_\bullet$ be a filtration on $P$ as in property (P).
The short exact sequence of Lemma \ref{lemma-property-P-sequence}
produces a distinguished triangle. Hence by
Derived Categories, Lemma \ref{derived-lemma-representable-homological}
it suffices to show that
$$
\Hom_{K(\text{Mod}_{(A, \text{d})})}(F_iP, N) = 0
$$
for all acyclic differential graded $A$-modules $N$ and all $i$.
Each of the differential graded modules $F_iP$ has a finite filtration
by admissible monomorphisms, whose graded pieces are direct sums
of shifts $A[k]$. Thus it suffices to prove that
$$
\Hom_{K(\text{Mod}_{(A, \text{d})})}(A[k], N) = 0
$$
for all acyclic differential graded $A$-modules $N$ and all $k$.
This follows from Lemma \ref{lemma-hom-from-shift-free}.
\end{proof}

\begin{lemma}
\label{lemma-right-resolution}
Let $(A, \text{d})$ be a differential graded algebra.
Let $M$ be a differential graded $A$-module. There exists a homomorphism
$M \to I$ of differential graded $A$-modules such that
\begin{enumerate}
\item $M \to I$ is a quasi-isomorphism, and
\item $I$ has property (I).
\end{enumerate}
\end{lemma}

\begin{proof}
Set $M = M_0$. We inductively choose short exact sequences
$$
0 \to M_i \to I_i \to M_{i + 1} \to 0
$$
where the maps $M_i \to I_i$ are chosen as in Lemma \ref{lemma-good-sub}.
This gives a ``resolution''
$$
0 \to M \to I_0 \xrightarrow{f_0} I_1 \xrightarrow{f_1} I_1 \to \ldots
$$
Denote $I$ the differential graded $A$-module with graded parts
$$
I^n = \prod\nolimits_{i \geq 0} I^{n - i}_i
$$
and differential defined by
$$
\text{d}_I(x) = f_i(x) + (-1)^i \text{d}_{I_i}(x)
$$
for $x \in I_i^{n - i}$. With these conventions $I$ is indeed a differential
graded $A$-module. Recalling that each $I_i$ has a two step filtration
$0 \to I_i' \to I_i \to I_i'' \to 0$ we set
$$
F_{2i}I^n = \prod\nolimits_{j \geq i} I^{n - j}_j
\subset
\prod\nolimits_{i \geq 0} I^{n - i}_i = I^n
$$
and we add a factor $I'_{i + 1}$ to $F_{2i}I$ to get $F_{2i + 1}I$.
These are differential graded submodules and the successive quotients
are products of shifts of $A^\vee$. By
Lemma \ref{lemma-source-graded-injective} we see that
the inclusions $F_{i + 1}I \to F_iI$ are admissible monomorphisms.
Finally, we have to show that the map $M \to I$ (given by the
augmentation $M \to I_0$) is a quasi-isomorphism. This follows from
Homology, Lemma \ref{homology-lemma-good-right-resolution-gives-qis}.
\end{proof}

\begin{lemma}
\label{lemma-hom-into-shift-dual-free}
Let $(A, \text{d})$ be a differential graded algebra. Then we have
$$
\Hom_{\text{Mod}_{(A, \text{d})}}(M, A^\vee[k]) =
\Ker(\text{d} : (M^\vee)^k \to (M^\vee)^{k + 1})
$$
and
$$
\Hom_{K(\text{Mod}_{(A, \text{d})})}(M, A^\vee[k]) = H^k(M^\vee)
$$
as functors in the differential graded $A$-module $M$.
\end{lemma}

\begin{proof}
This is clear from the discussion above.
\end{proof}

\begin{lemma}
\label{lemma-property-I-sequence}
Let $(A, \text{d})$ be a differential graded algebra.
Let $I$ be a differential graded $A$-module. If $F_\bullet$
is a filtration as in property (I), then we obtain an
admissible short exact sequence
$$
0 \to I \to
\prod\nolimits I/F_iI \to
\prod\nolimits I/F_iI \to 0
$$
of differential graded $A$-modules.
\end{lemma}

\begin{proof}
Omitted. Hint: This is dual to Lemma \ref{lemma-property-P-sequence}.
\end{proof}

\begin{proposition}
\label{proposition-homotopy-category-triangulated}
Let $(A, \text{d})$ be a differential graded algebra. The homotopy category
$K(\text{Mod}_{(A, \text{d})})$ of differential graded $A$-modules with its
natural translation functors and distinguished triangles is a triangulated
category.
\end{proposition}

\begin{proof}
We know that $K(\text{Mod}_{(A, \text{d})})$ is a pre-triangulated category.
Hence it suffices to prove TR4 and to prove it we can use
Derived Categories, Lemma \ref{derived-lemma-easier-axiom-four}.
Let $K \to L$ and $L \to M$ be composable morphisms of
$K(\text{Mod}_{(A, \text{d})})$. By
Lemma \ref{lemma-sequence-maps-split} we may assume that
$K \to L$ and $L \to M$ are admissible monomorphisms.
In this case the result follows from
Lemma \ref{lemma-two-split-injections}.
\end{proof}

\begin{lemma}
\label{lemma-hom-from-shift-free}
Let $(A, d)$ be a differential graded algebra. Then we have
$$
\Hom_{\text{Mod}_{(A, \text{d})}}(A[k], M) =
\Ker(\text{d} : M^{-k} \to M^{-k + 1})
$$
and
$$
\Hom_{K(\text{Mod}_{(A, \text{d})})}(A[k], M) = H^{-k}(M)
$$
for any differential graded $A$-module $M$.
\end{lemma}

\begin{proof}
Immediate from the definitions.
\end{proof}

\begin{lemma}
\label{lemma-good-sub}
Let $(A, \text{d})$ be a differential graded algebra.
Let $M$ be a differential graded $A$-module. There exists a homomorphism
$M \to I$ of differential graded $A$-modules with the following
properties
\begin{enumerate}
\item $M \to I$ is injective,
\item $\Coker(\text{d}_M) \to \Coker(\text{d}_I)$ is injective,
and
\item $I$ sits in an admissible short exact sequence
$0 \to I' \to I \to I'' \to 0$ where $I'$, $I''$ are products
of shifts of $A^\vee$.
\end{enumerate}
\end{lemma}

\begin{proof}
We will use the functors $N \mapsto N^\vee$ (from left to right
differential graded modules and from right to left differential
graded modules) constructed in
Section \ref{section-modules-noncommutative-differential-graded}
and all of their properties.
For every $k \in \mathbf{Z}$ let $Q_k$ be the free left $A$-module with
generators $x, y$ in degrees $k$ and $k + 1$. Define the structure of a
left differential graded $A$-module on $Q_k$ by setting $\text{d}(x) = y$
and $\text{d}(y) = 0$. Arguing exactly as in the proof of
Lemma \ref{lemma-good-quotient} we find a surjection
$$
\bigoplus\nolimits_{i \in I} Q_{k_i} \longrightarrow M^\vee
$$
of left differential graded $A$-modules. Then we can consider the injection
$$
M \to (M^\vee)^\vee \to (\bigoplus\nolimits_{i \in I} Q_{k_i})^\vee =
\prod\nolimits_{i \in I} I_{k_i}
$$
where $I_k = Q_{-k}^\vee$ is the ``dual'' right differential graded $A$-module.
Further, the short exact sequence $0 \to A[-k - 1] \to Q_k \to A[-k] \to 0$
produces a short exact sequence
$0 \to A^\vee[k] \to I_k \to A^\vee[k + 1] \to 0$.

\medskip\noindent
The result of the previous paragraph produces $M \to I$
having properties (1) and (3). To obtain property (2), suppose
$\overline{m} \in \Coker(\text{d}_M)$ is a nonzero element of
degree $k$. Pick a map $\lambda : M^k \to \mathbf{Q}/\mathbf{Z}$
which vanishes on $\Im(M^{k - 1} \to M^k)$ but not on $m$. By
Lemma \ref{lemma-hom-into-shift-dual-free}
this corresponds to a homomorphism $M \to A^\vee[k]$ of
differential graded $A$-modules which does not vanish on $m$.
Hence we can achieve (2) by adding
a product of copies of shifts of $A^\vee$.
\end{proof}

\begin{lemma}
\label{lemma-source-graded-injective}
Let $(A, \text{d})$ be a differential graded algebra.
Let $I \to M$ be an injective homomorphism of differential graded
$A$-modules. If $I$ is graded injective, then
$I \to M$ is an admissible monomorphism.
\end{lemma}

\begin{proof}
This is immediate from the definitions.
\end{proof}

\begin{lemma}
\label{lemma-sequence-maps-split}
Let $(A, \text{d})$ be a differential graded algebra.
Let $L_1 \to L_2 \to \ldots \to L_n$
be a sequence of composable homomorphisms of
differential graded $A$-modules.
There exists a commutative diagram
$$
\xymatrix{
L_1 \ar[r] &
L_2 \ar[r] &
\ldots \ar[r] &
L_n \\
M_1 \ar[r] \ar[u] &
M_2 \ar[r] \ar[u] &
\ldots \ar[r] &
M_n \ar[u]
}
$$
in $\text{Mod}_{(A, \text{d})}$ such that each $M_i \to M_{i + 1}$
is an admissible monomorphism and each $M_i \to L_i$
is a homotopy equivalence.
\end{lemma}

\begin{proof}
The case $n = 1$ is without content.
Lemma \ref{lemma-make-injective} is the case $n = 2$.
Suppose we have constructed the diagram
except for $M_n$. Apply Lemma \ref{lemma-make-injective} to
the composition $M_{n - 1} \to L_{n - 1} \to L_n$.
The result is a factorization $M_{n - 1} \to M_n \to L_n$
as desired.
\end{proof}

\begin{lemma}
\label{lemma-two-split-injections}
Let $(A, \text{d})$ be a differential graded algebra. Suppose that
$\alpha : K \to L$ and $\beta : L \to M$ are admissible monomorphisms
of differential graded $A$-modules. Then there exist distinguished triangles
$(K, L, Q_1, \alpha, p_1, d_1)$, $(K, M, Q_2, \beta \circ \alpha, p_2, d_2)$
and $(L, M, Q_3, \beta, p_3, d_3)$ for which TR4 holds.
\end{lemma}

\begin{proof}
Say $\pi_1 : L \to K$ and $\pi_3 : M \to L$ are homomorphisms
of graded $A$-modules which are left inverse to $\alpha$ and $\beta$.
Then also $K \to M$ is an admissible monomorphism with left
inverse $\pi_2 = \pi_1 \circ \pi_3$. Let us write $Q_1$, $Q_2$
and $Q_3$ for the cokernels of $K \to L$, $K \to M$, and $L \to M$.
Then we obtain identifications (as graded $A$-modules)
$Q_1 = \Ker(\pi_1)$, $Q_3 = \Ker(\pi_3)$ and
$Q_2 = \Ker(\pi_2)$. Then $L = K \oplus Q_1$ and
$M = L \oplus Q_3$ as graded $A$-modules. This implies
$M = K \oplus Q_1 \oplus Q_3$. Note that $\pi_2 = \pi_1 \circ \pi_3$
is zero on both $Q_1$ and $Q_3$. Hence $Q_2 = Q_1 \oplus Q_3$.
Consider the commutative diagram
$$
\begin{matrix}
0 & \to & K & \to & L & \to & Q_1 & \to & 0 \\
  &     & \downarrow &     & \downarrow &     & \downarrow  & \\
0 & \to & K & \to & M & \to & Q_2 & \to & 0 \\
  &     & \downarrow &     & \downarrow &     & \downarrow  & \\
0 & \to & L & \to & M & \to & Q_3 & \to & 0
\end{matrix}
$$
The rows of this diagram are admissible short exact sequences, and
hence determine distinguished triangles by definition. Moreover
downward arrows in the diagram above are compatible with the chosen
splittings and hence define morphisms of triangles
$$
(K \to L \to Q_1 \to K[1])
\longrightarrow
(K \to M \to Q_2 \to K[1])
$$
and
$$
(K \to M \to Q_2 \to K[1])
\longrightarrow
(L \to M \to Q_3 \to L[1]).
$$
Note that the splittings $Q_3 \to M$ of the bottom sequence in the
diagram provides a splitting for the split sequence
$0 \to Q_1 \to Q_2 \to Q_3 \to 0$ upon composing with $M \to Q_2$.
It follows easily from this that the morphism $\delta : Q_3 \to Q_1[1]$
in the corresponding distinguished triangle
$$
(Q_1 \to Q_2 \to Q_3 \to Q_1[1])
$$
is equal to the composition $Q_3 \to L[1] \to Q_1[1]$.
Hence we get a structure as in the conclusion of axiom TR4.
\end{proof}

\section{Injective modules and differential graded algebras}
\label{section-modules-noncommutative-differential-graded}

\begin{lemma}
\label{lemma-make-injective}
Let $(A, \text{d})$ be a differential graded algebra.
Let $\alpha : K \to L$ be a homomorphism of differential graded
$A$-modules. There exists a factorization
$$
\xymatrix{
K \ar[r]^{\tilde \alpha} \ar@/_1pc/[rr]_\alpha &
\tilde L \ar[r]^\pi & L
}
$$
in $\text{Mod}_{(A, \text{d})}$ such that
\begin{enumerate}
\item $\tilde \alpha$ is an admissible monomorphism (see
Definition \ref{definition-admissible-ses}),
\item there is a morphism $s : L \to \tilde L$
such that $\pi \circ s = \text{id}_L$ and such that
$s \circ \pi$ is homotopic to $\text{id}_{\tilde L}$.
\end{enumerate}
\end{lemma}

\begin{proof}
The proof is identical to the proof of
Derived Categories, Lemma \ref{derived-lemma-make-injective}.
Namely, we set $\tilde L = L \oplus C(1_K)$ and we use elementary
properties of the cone construction.
\end{proof}

\begin{definition}
\label{definition-admissible-ses}
Let $(A, \text{d})$ be a differential graded algebra.
\begin{enumerate}
\item A homomorphism $K \to L$ of differential graded $A$-modules
is an {\it admissible monomorphism} if there exists a graded $A$-module
map $L \to K$ which is left inverse to $K \to L$.
\item A homomorphism $L \to M$ of differential graded $A$-modules
is an {\it admissible epimorphism} if there exists a graded $A$-module
map $M \to L$ which is right inverse to $L \to M$.
\item A short exact sequence $0 \to K \to L \to M \to 0$ of differential
graded $A$-modules is an {\it admissible short exact sequence}
if it is split as a sequence of graded $A$-modules.
\end{enumerate}
\end{definition}

\begin{lemma}
\label{lemma-compose-tensor-functors-general-algebra}
Let $R$ be a ring. Let $(A, \text{d})$, $(B, \text{d})$, and
$(C, \text{d})$ be differential graded $R$-algebras. Assume
that (\ref{equation-plain-versus-derived-algebras}) is an isomorphism.
Let $N$ be a differential graded $(A, B)$-bimodule.
Let $N'$ be a differential graded $(B, C)$-bimodule.
Then the composition
$$
\xymatrix{
D(A, \text{d}) \ar[rr]^{- \otimes_A^\mathbf{L} N} & &
D(B, \text{d}) \ar[rr]^{- \otimes_B^\mathbf{L} N'} & &
D(C, \text{d})
}
$$
is isomorphic to $- \otimes_A^\mathbf{L} N''$ for a differential graded
$(A, C)$-bimodule $N''$ described in the proof.
\end{lemma}

\begin{proof}
By Lemma \ref{lemma-functoriality-bc} we may replace $N$ and $N'$ by
quasi-isomorphic bimodules. Thus we may assume $N$, resp.\ $N'$
has property (P) as differential graded
$(A, B)$-bimodule, resp.\ $(B, C)$-bimodule, see
Lemma \ref{lemma-bimodule-resolve}. We claim the lemma holds
with the $(A, C)$-bimodule $N'' = N \otimes_B N'$.
To prove this, it suffices to show that
$$
N_B \otimes_B^\mathbf{L} N' \longrightarrow (N \otimes_B N')_C
$$
is an isomorphism in $D(C, \text{d})$, see
Lemma \ref{lemma-compose-tensor-functors-general}.

\medskip\noindent
Let $F_\bullet$ be the filtration on $N$ as in property (P) for bimodules.
By Lemma \ref{lemma-bimodule-property-P-sequence}
there is a short exact sequence
$$
0 \to
\bigoplus\nolimits F_iN \to
\bigoplus\nolimits F_iN \to N \to 0
$$
of differential graded $(A, B)$-bimodules which is split as a sequence
of graded $(A, B)$-bimodules. A fortiori this is an admissible short exact
sequence of differential graded $B$-modules and this produces a distinguished
triangle
$$
\bigoplus\nolimits F_iN_B \to
\bigoplus\nolimits F_iN_B \to N_B \to
\bigoplus\nolimits F_iN_B[1]
$$
in $D(B, \text{d})$. Using that $- \otimes_B^\mathbf{L} N'$
is an exact functor of triangulated categories and commutes
with direct sums and using that $- \otimes_B N'$ transforms
admissible exact sequences into admissible exact sequences
and commutes with direct sums we reduce to proving
that
$$
(F_pN)_B \otimes_B^\mathbf{L} N' \longrightarrow (F_pN)_B \otimes_B N'
$$
is a quasi-isomorphism for all $p$. Repeating the argument
with the short exact sequences of $(A, B)$-bimodules
$$
0 \to F_pN \to F_{p + 1}N \to F_{p + 1}N/F_pN \to 0
$$
which are split as graded $(A, B)$-bimodules
we reduce to showing the same statement for $F_{p + 1}N/F_pN$.
Since these modules are direct sums of shifts of $(A \otimes_R B)_B$
we reduce to showing that
$$
(A \otimes_R B)_B \otimes_B^\mathbf{L} N'
\longrightarrow
(A \otimes_R B)_B \otimes_B N'
$$
is a quasi-isomorphism.

\medskip\noindent
Choose a filtration $F_\bullet$ on $N'$ as in property (P) for bimodules.
Choose a quasi-isomorphism $P \to (A \otimes_R B)_B$
of differential graded $B$-modules where $P$ has property (P).
We have to show that
$P \otimes_B N' \to (A \otimes_R B)_B \otimes_B N'$ is
a quasi-isomorphism because $P \otimes_B N'$ represents
$(A \otimes_R B)_B \otimes_B^\mathbf{L} N'$ in $D(C, \text{d})$
by the construction in Lemma \ref{lemma-derived-bc}.
As $N' = \colim F_pN'$ we find
that it suffices to show that
$P \otimes_B F_pN' \to (A \otimes_R B)_B \otimes_B F_pN'$
is a quasi-isomorphism. Using the short exact sequences
$0 \to F_pN' \to F_{p + 1}N' \to F_{p + 1}N'/F_pN' \to 0$
which are split as graded $(B, C)$-bimodules we reduce to showing
$P \otimes_B F_{p + 1}N'/F_pN' \to
(A \otimes_R B)_B \otimes_B F_{p + 1}N'/F_pN'$
is a quasi-isomorphism for all $p$.
Then finally using that $F_{p + 1}N'/F_pN'$
is a direct sum of shifts of ${}_B(B \otimes_R C)_C$
we conclude that it suffices to show that
$$
P \otimes_B {}_B(B \otimes_R C)_C \to
(A \otimes_R B)_B \otimes_B {}_B(B \otimes_R C)_C
$$
is a quasi-isomorphism. Since $P \to (A \otimes_R B)_B$
is a resolution by a module satisfying property (P)
this map of differential graded $C$-modules
represents the morphism (\ref{equation-plain-versus-derived-algebras})
in $D(C, \text{d})$ and the proof is complete.
\end{proof}

\end{document}
